\section{Introduction}

Time-series forecasting plays a central role in many scientific and operational applications, including weather prediction, energy management, transportation, finance, and healthcare. In such settings, an accurate point forecast alone is often insufficient: decision makers must also understand the uncertainty associated with each prediction. For example, in weather forecasting, the consequences of underestimating uncertainty may range from inefficient resource allocation to inadequate preparation for extreme events. Reliable prediction intervals are therefore essential for communicating forecast uncertainty and supporting risk-aware decision making.

Traditional approaches to uncertainty quantification typically rely on parametric assumptions about the forecast errors or on probabilistic models that explicitly estimate the conditional distribution of future observations. However, these assumptions may be difficult to verify and can lead to poorly calibrated uncertainty estimates when the forecasting model is misspecified or when the data distribution changes over time. These limitations are particularly important in modern time-series applications, where flexible machine-learning and deep-learning models can achieve high predictive accuracy but do not naturally provide prediction intervals with reliable statistical coverage.

Conformal Prediction (CP) \cite{vovk2005conformal} provides a model-agnostic framework for converting the output of an arbitrary prediction model into prediction intervals with a user-specified coverage level. Under exchangeability, conformal methods offer finite-sample, distribution-free marginal coverage guarantees without requiring a correct probabilistic model for the data. This combination of flexibility and statistical validity has made conformal prediction an attractive approach to uncertainty quantification in classification and regression tasks. Its application to time-series forecasting, however, is not straightforward because observations collected over time are generally dependent and nonstationary, violating the exchangeability assumption underlying classical conformal guarantees.



Several studies have extended conformal prediction to time-series forecasting, where temporal dependence and distribution shifts violate the exchangeability assumption used by standard conformal methods.  A simple approach is to use a rolling window mechanism where the calibration set is dynamically updated. At every time step, the newest nonconformity score is added and the oldest is removed. This handles nonstationarity by assuming that recent scores are more relevant than old scores. Xu and Xie proposed EnbPI, which constructs sequential prediction intervals using bootstrap ensemble forecasts and a sliding window of recent residuals \cite{xu2021enbpi}. They later introduced SPCI, which models the conditional distribution of future nonconformity scores based on previous scores, allowing the interval width to adapt to temporal changes in forecast uncertainty \cite{xu2023spci}. Gibbs and Candes proposed adaptive conformal inference, which dynamically updates the target miscoverage level and provides long-run coverage control under arbitrary distribution shifts \cite{gibbs2021adaptive}.
 Stankevičiūtė et al. introduced conformal time-series forecasting for recurrent neural networks, calibrating prediction regions across multiple forecast horizons \cite{stankeviciute2021conformal}. Auer et al. proposed HopCPT, which uses modern Hopfield networks to select historical nonconformity scores that are relevant to the current forecasting context \cite{auer2023hopcpt}.
% These methods improve conformal prediction for dependent data through different forms of adaptivity and validity, depending on their assumptions regarding temporal dependence and distribution shift.
 Despite these advances, existing time-series conformal methods primarily adapt to changes in the residual distribution over time. They generally do not exploit model-based, input-dependent uncertainty to adapt interval width to the difficulty of the particular forecast instance. Consequently, temporal adaptation and instance-dependent adaptation have largely been treated separately.

%In this work, we study instance-dependent conformal prediction for time-series forecasting, with particular emphasis on weather data. We consider probabilistic Mixture-of-Expert (MoE) models that provide both a point prediction and an input-dependent estimate of predictive uncertainty. The uncertainty estimate is used to normalize the forecast error, and conformal calibration is then applied to obtain prediction intervals whose widths adapt to the current forecasting conditions. To address temporal nonstationarity, the calibration scores are maintained using a rolling calibration window. The resulting framework combines two complementary forms of adaptivity: instance-level adaptation through the predicted uncertainty and temporal adaptation through continuously updated conformal calibration.  Weather forecasting provides a particularly challenging setting for evaluating this approach because weather observations exhibit strong temporal dependence, seasonality, spatial heterogeneity, heteroscedastic errors, and occasional abrupt changes. Our experiments examine whether instance-dependent scaling can produce narrower and more informative prediction intervals than fixed-width conformal methods while maintaining the desired empirical coverage over time. The proposed framework is model-agnostic and can be applied on top of a broad range of neural time-series forecasting architectures.


In this work, we study instance-dependent conformal prediction for  time-series forecasting. 
We consider probabilistic forecasting models that provide both a point forecast and an input-dependent estimate of predictive uncertainty. The predicted uncertainty is used to normalize the forecast error, producing conformal prediction intervals whose widths adapt to the difficulty of each individual forecast.



Our main contribution is a conformal forecasting framework that combines two complementary forms of adaptation. First, instance-dependent predictive uncertainty is used to normalize nonconformity scores, allowing interval widths to reflect the predicted difficulty of individual forecasts. Second, a rolling calibration window with adaptive conformal inference continuously updates the conformal threshold to account for temporal changes in the residual distribution. The resulting interval width therefore adapts both across instances and over time. We instantiate the framework using Gaussian and mixture-of-Gaussian forecasting models and show that the combined adaptation yields substantially narrower intervals than CQR \cite{cqr2019} while maintaining empirical coverage close to the target level.